% SHARD-ID: AQM-52-F3-REGULAR-FUNCTION-SYMMETRY-LINE-BY-LINE
% SHARD-TITLE: The \(\mathbb F_3\) Regular-Function Symmetry Action Line by Line
% SHARD-SUMMARY: Computes \(g_\#F=F\circ g^{-1}\) for all six affine symmetries of \(\mathbf A^1_{\mathbb F_3}\).
% SHARD-SUMMARY: Checks residue-profile covariance at rational and quadratic closed points.
% SHARD-SUMMARY: Computes the transformed reduced Weyl smearings and the repeated-factor Artin support transport.
% SHARD-KEYWORDS: F3, affine symmetry, regular function, residue profile, Weyl smearing, Artin ring

\section{The \texorpdfstring{\(\mathbb F_3\)}{F3} Regular-Function Symmetry Action Line by Line}
\label{sec:f3-regular-function-symmetry-line-by-line}

\subsection{Status and Source Anchors}
This is the \(k=\mathbb F_3\) companion to AQM-51.  It uses AQM-46 for the
affine symmetry group, AQM-48 for the full closed-point Hilbert action,
AQM-49 for regular functions and reduced smearing, and AQM-50 for the concrete
regular-function examples.

\subsection{Lamport Dependency Skeleton}
\[
\begin{array}{rcl}
D1&:&R=\mathbb F_3[x],\ X=\Spec R,\ G=\{g_{a,b}:x\mapsto ax+b\}.\\
D2&:&g_\#F=F(a^{-1}(x-b)).\\
L1&:&g_\#x=a^{-1}(x-b)\text{ for all six }g.\\
D3&:&\tau:x\mapsto x+1,\quad \rho:x\mapsto2x.\\
T1&:&\operatorname{ev}(g_\#F)=S_g^\Pi\operatorname{ev}(F)\text{ at rational points.}\\
D4&:&h=x(x-1),\quad F=1,\quad G=x+1.\\
T2&:&\alpha_\tau W_h^{\rm red}(F,G)=W_{\tau_\#h}^{\rm red}(\tau_\#F,\tau_\#G).\\
D5&:&q_0=x^2+1,\quad q_1=x^2+x+2.\\
T3&:&\tau_\#q_0=q_1\text{ and quadratic labels transform by }\tau_{g,\pi}^{-1}.\\
D6&:&m=x^2q_0^3.\\
T4&:&\tau_\#m=(x-1)^2q_1^3,\text{ preserving reduced and thickened dimensions.}
\end{array}
\]

\subsection{Step 1: The Six Symmetries D1--L1}
For \(k=\mathbb F_3\),
\[
  G=\operatorname{Aut}_{\mathbb F_3}(X)=
  \{g_{a,b}:x\mapsto ax+b,\ a\in\{1,2\},\ b\in\mathbb F_3\}.
\]
By AQM-51,
\[
  g_\#F=F\circ g^{-1},\qquad g_\#F(x)=F(a^{-1}(x-b)).
\]
Since \(1^{-1}=1\) and \(2^{-1}=2\) in \(\mathbb F_3\),
\[
\begin{array}{c|c|c}
g&(a,b)&g_\#x\\
\hline
x&(1,0)&x\\
x+1&(1,1)&x-1=x+2\\
x+2&(1,2)&x-2=x+1\\
2x&(2,0)&2x\\
2x+1&(2,1)&2(x-1)=2x+1\\
2x+2&(2,2)&2(x-2)=2x+2
\end{array}
\]

\paragraph{Lemma \(L1\).}
The table gives \(g_\#x=g^{-1}(x)\) in every case.

\emph{Proof.}
The inverse of \(g_{a,b}\) is \(g_{a^{-1},-a^{-1}b}\).  Therefore
\(g_\#x=a^{-1}(x-b)\).  Substituting the six pairs \((a,b)\) gives the table.

\subsection{Step 2: Translation and Reflection D3}
Write
\[
  \tau:x\mapsto x+1,\qquad \rho:x\mapsto2x.
\]
Then \(\tau_\#F=F(x-1)\) and \(\rho_\#F=F(2x)\).  For \(F=x\),
\[
  \tau_\#x=x-1,\qquad \rho_\#x=2x.
\]
For the local regular function \(x^{-1}\in\Gamma(D(x),\mathcal O_X)\),
\[
  \tau_\#(x^{-1})=(x-1)^{-1}\in\Gamma(D(x-1),\mathcal O_X),
  \rho_\#(x^{-1})=(2x)^{-1}=2x^{-1}\in\Gamma(D(x),\mathcal O_X).
\]
The open \(D(x)\) moves to \(D(x-1)\) under \(\tau\), because \(x_0\) moves to
\(x_1\).  It is fixed by \(\rho\), because \(x_0\) is fixed by \(\rho\).

\subsection{Step 3: Rational-Point Profile Check T1}
The rational closed points are \(x_r=(x-r)\), \(r=0,1,2\).  The translation
\(\tau\) sends \(x_r\mapsto x_{r+1}\).
At rational points the residue fields are all \(\mathbb F_3\), and the
residue-field isomorphism is the identity after evaluation.

Check \(F=x\).  The source value at \(x_r\) is \(r\).  The transformed function
is \(\tau_\#F=x-1\).  At \(x_{r+1}\), \((\tau_\#F)(x_{r+1})=(r+1)-1=r\).
Line by line,
\[
\begin{array}{c|c|c|c}
\text{source}&F&\text{target}&\tau_\#F\\
\hline
x_0&0&x_1&1-1=0\\
x_1&1&x_2&2-1=1\\
x_2&2&x_0&0-1=2
\end{array}
\]
This is exactly \(\operatorname{ev}(\tau_\#F)=S_\tau^\Pi\operatorname{ev}(F)\).

For \(x^{-1}\) on \(D(x)\), the source points are \(x_1,x_2\).  They move to
\(x_2,x_0\).  The values are
\[
\begin{array}{c|c|c|c}
\text{source}&x^{-1}&\text{target}&(x-1)^{-1}\\
\hline
x_1&1&x_2&(2-1)^{-1}=1\\
x_2&2&x_0&(0-1)^{-1}=2
\end{array}
\]
Thus local regular functions satisfy the same covariance rule.

\subsection{Step 4: A Reduced Two-Qutrit Smearing D4--T2}
Use the AQM-50 support and labels
\[
  h=x(x-1),\qquad F=1,\qquad G=x+1.
\]
The support is \(Z_{\rm cl}(h)=\{x_0,x_1\}\).  The source labels are
\(x_0:(1,1)\) and \(x_1:(1,2)\).
Under \(\tau\),
\[
  \tau_\#h=(x-1)(x-2),\qquad \tau_\#F=1,\qquad \tau_\#G=x.
\]
The target support is \(Z_{\rm cl}(\tau_\#h)=\{x_1,x_2\}\).  At \(x_1\) and
\(x_2\), the target labels are \((1,1)\) and \((1,2)\).

\paragraph{Theorem \(T2\).}
\[
  \alpha_\tau W_{x(x-1)}^{\rm red}(1,x+1)=
  W_{(x-1)(x-2)}^{\rm red}(1,x).
\]

\emph{Proof.}
The source label \(((1,1),(1,2))\) at \((x_0,x_1)\) has just been carried to
the same ordered values at \((x_1,x_2)\).  AQM-46 says
\(\alpha_\tau W(e)=W(S_\tau e)\), and AQM-51 identifies
\(S_\tau e\) with the reduced label of \((1,x)\) on the transformed support.

\subsection{Step 5: Quadratic Closed Point D5--T3}
Let \(q_0=x^2+1\) and \(q_1=x^2+x+2\).
For \(\tau\),
\[
  \tau_\#q_0=q_0(x-1)=(x-1)^2+1=x^2+x+2=q_1.
\]
Let \(K_{q_0}=\mathbb F_3[\alpha]\), \(\alpha^2=2\), and
\(K_{q_1}=\mathbb F_3[\beta]\), \(\beta^2+\beta+2=0\).
The residue pullback is
\[
  \tau_{\tau,q_0}:K_{q_1}\to K_{q_0},\qquad
  \beta\longmapsto\alpha+1.
\]
Hence \(\tau_{\tau,q_0}^{-1}(\alpha)=\beta-1=\beta+2\).

Use the AQM-50 labels
\[
  F=x+1,\qquad G=2x
\]
on \(q_0\).  The source label is \((\alpha+1,\ 2\alpha)\).  The transformed
functions are \(\tau_\#F=x\) and \(\tau_\#G=2x+1\), giving
\((\beta,\ 2\beta+1)\) at \(q_1\).

\paragraph{Theorem \(T3\).}
\[
  (\beta,\ 2\beta+1)=
  \bigl(\tau_{\tau,q_0}^{-1}(\alpha+1),
        \tau_{\tau,q_0}^{-1}(2\alpha)\bigr).
\]

\emph{Proof.}
Since \(\tau_{\tau,q_0}^{-1}(\alpha)=\beta+2\),
\[
  \tau_{\tau,q_0}^{-1}(\alpha+1)=\beta,\qquad
  \tau_{\tau,q_0}^{-1}(2\alpha)=2(\beta+2)=2\beta+1.
\]
These are exactly the target residues of \(x\) and \(2x+1\) at \(q_1\).

\subsection{Step 6: Repeated Factors D6--T4}
Take the AQM-50 repeated-factor polynomial \(m=x^2q_0^3\).
Under \(\tau\),
\[
  \tau_\#m=(x-1)^2q_1^3.
\]
The reduced source support is \(\{x_0,x_{q_0}\}\), and the reduced target
support is \(\{x_1,x_{q_1}\}\).  The reduced Hilbert dimension is unchanged:
\[
  3\cdot9=27.
\]
The thickened source rings are \(\mathbb F_3[x]/(x^2)\) and
\(\mathbb F_3[x]/(q_0^3)\).  The target rings are
\(\mathbb F_3[x]/((x-1)^2)\) and \(\mathbb F_3[x]/(q_1^3)\).  The map
\(F\mapsto F(x-1)\) gives ring isomorphisms from source to target.  Therefore
the thickened dimension is also unchanged:
\[
  3^2\cdot3^6=3^8=6561.
\]

\paragraph{Theorem \(T4\).}
The affine action preserves the exponents \(2\) and \(3\).  It moves the
closed support and the thickened Artin rings, but it does not turn a repeated
factor into extra closed points.

\emph{Proof.}
The equality \(\tau_\#m=(x-1)^2q_1^3\) shows that exponents are carried with
their irreducible factors.  The reduced support keeps only the distinct
factors, giving one \(3\)-level site and one \(9\)-level site before and after
translation.  The thickened rings keep the powers, so their sizes remain
\(3^2\) and \(3^6\).  Thus multiplicity is preserved as Artin-ring thickness,
not converted into a larger closed-point set.

\subsection{Step 7: The Full Hilbert-Space Statement}
The full closed-point Hilbert representation from AQM-48 is
\[
  \mathcal H^0=\ell^2\!\left(\bigoplus_{\pi}K_\pi\right),
\]
where \(\pi\) runs over all monic irreducibles and vectors have finite closed
support.  The affine unitary satisfies
\[
  U_gW_h^{\rm red}(F,G)U_g^*=W_{g_\#h}^{\rm red}(g_\#F,g_\#G)
\]
for every finite support selector \(h\) and polynomial labels \(F,G\).

There is no operator \(W(\operatorname{ev}(x))\) in the algebraic
finite-support Weyl net, because AQM-50 proves that \(x\) has infinite closed
support.  The symmetry acts on that infinite profile by relabelling it, but
the Weyl operator exists only after choosing a finite support such as \(h\), or
after adding a new infinite-volume completion convention.
